\subsection*{Finding an unnamed participant in a pictured passage}
The same bronze image can be easy to copy and difficult to choose. We asked who had a mirror and who lacked it in a passage that supplied the two names as bronze inscriptions (Fig.~\ref{fig:inverse-role-writer}c). The inverse-trained writer returned Qin for both questions. Across eight objects it gave 17/32 correct names and followed the correct participant through both passage orders in 7/16 pairs, below the prespecified target of 12. The randomized control gave 15/32 and 6/16. The matched writer returned the first participant in 27 of the 32 responses. Although it could write unfamiliar ink, this positional preference frequently displaced the requested participant. With four historical participants and distinct objects it gave 5/96 correct answers, compared with 13 under randomization. Neither writer followed any object through all four rotations. Familiar reordered descriptions remained relatively strong (344/404), whereas new layouts fell from 297 to 11/404.

On human-written traditional Chinese questions asking for an unnamed person, the matched writer improved from 4 to 62/387 exact answers after adaptation. Both arms received the same ordinary DRCD supervision. Randomizing the additional competing assignments gave 64/387, leaving a matched difference of $-0.5$ percentage points (passage-bootstrap 95\% interval $-3.1$ to $2.1$). Each fitted arm answered both questions correctly on just 1/131 same-passage pairs. The separate passage reader gave 95/387 and 7/131. The matched writer gave 549/3,505 complete general DRCD responses and zero of 571 canonical DuReader answers. These experiments used 694,152 network parameters with inherited weights and a fixed writing memory. The two final adaptations took 95.8 and 91.2 seconds on one RTX 4090 D. These costs exclude the inherited training.

The saved pixels identify two sources of error. Of 325 incorrect human who-answers, 314 traced to an uninterrupted span of the supplied passage, and 301 differed from every reference at the first glyph. Figure~\ref{fig:inverse-role-writer}d shows a place name copied in response to a question asking for an inventor. Other errors began correctly but continued past the name, as when the writer added a verb fragment after William Jones. All 387 questions admitted an answer-containing sentence and a sentence without an annotated answer. Supplying the shortest sentence containing a reference raised complete answers from 62 to 141; the length-matched distractors yielded none. A selector available at inference, based on overlap between question and sentence glyphs, retained an annotated answer for 301 questions and raised exact responses to 116/387 (gain 14.0 percentage points, passage-bootstrap 95\% interval 9.1--18.8). Applying exactly the same sentence selector gave 127 correct answers for the randomized writer and 100 for the passage reader. On the 131 paired questions, sentence selection yielded nine both-correct pairs for matched training, 12 under randomization and eight for the passage reader. The gain from narrowing the evidence therefore transfers across weight sets. Even with an annotated evidence sentence, 246 answers remained incorrect. In 27 human question pairs, one sentence contained different valid answers to both questions. With overlap-based sentence selection the matched writer answered both correctly in two pairs, compared with four under randomization and five for the passage reader.
